\documentclass[11pt,a4paper]{article}

\usepackage[T1]{fontenc}
\usepackage{mathtools,amssymb,amsthm}
\usepackage{newtxtext,newtxmath}
\usepackage{microtype}
\usepackage[a4paper,margin=25mm,headheight=14pt]{geometry}
\usepackage{booktabs,longtable,array,tabularx}
\usepackage{enumitem}
\usepackage{xcolor}
\usepackage[most]{tcolorbox}
\usepackage{listings}
\usepackage{fancyhdr}
\usepackage{hyperref}
\usepackage[nameinlink,capitalise,noabbrev]{cleveref}
\usepackage{url}

\definecolor{deepblue}{RGB}{25,55,95}
\definecolor{softblue}{RGB}{238,244,251}
\definecolor{softgray}{RGB}{246,246,246}
\definecolor{deepred}{RGB}{125,35,35}
\hypersetup{
  colorlinks=true,
  linkcolor=deepblue,
  citecolor=deepblue,
  urlcolor=deepred,
  pdfauthor={Research article prepared for the ProveIt/Fabius Function corpus},
  pdftitle={Rademacher Transseries for the Partition Numbers and Their Inverse},
  pdfsubject={Partition numbers, Rademacher series, transseries, inverse asymptotics, Lambert W, Bell polynomials}
}

\pagestyle{fancy}
\fancyhf{}
\fancyhead[L]{\small Rademacher transseries for partition numbers}
\fancyhead[R]{\small \thepage}
\renewcommand{\headrulewidth}{0.4pt}

\setlength{\parindent}{1.25em}
\setlength{\parskip}{0.15em}
\setlist{nosep,leftmargin=2.1em}
\allowdisplaybreaks
\numberwithin{equation}{section}

\newtheorem{theorem}{Theorem}[section]
\newtheorem{proposition}[theorem]{Proposition}
\newtheorem{lemma}[theorem]{Lemma}
\newtheorem{corollary}[theorem]{Corollary}
\theoremstyle{definition}
\newtheorem{definition}[theorem]{Definition}
\newtheorem{algorithm}[theorem]{Algorithm}
\theoremstyle{remark}
\newtheorem{remark}[theorem]{Remark}
\newtheorem{warning}[theorem]{Warning}

\newtcolorbox{resultbox}{
  enhanced, breakable, colback=softblue, colframe=deepblue,
  boxrule=0.7pt, arc=1.2mm, left=2mm,right=2mm,top=1.5mm,bottom=1.5mm}
\newtcolorbox{cautionbox}{
  enhanced, breakable, colback=softgray, colframe=deepred,
  boxrule=0.7pt, arc=1.2mm, left=2mm,right=2mm,top=1.5mm,bottom=1.5mm}

\lstset{
  basicstyle=\ttfamily\small,
  columns=fullflexible,
  breaklines=true,
  frame=single,
  backgroundcolor=\color{softgray},
  rulecolor=\color{black!30},
  showstringspaces=false,
  keepspaces=true
}

\DeclareMathOperator{\Res}{Res}
\DeclareMathOperator{\ord}{ord}
\DeclareMathOperator{\sgn}{sgn}
\DeclareMathOperator{\Li}{Li}
\newcommand{\e}{\mathrm e}
\newcommand{\ii}{\mathrm i}
\newcommand{\dd}{\,\mathrm d}
\newcommand{\N}{\mathbb N}
\newcommand{\Z}{\mathbb Z}
\newcommand{\R}{\mathbb R}
\newcommand{\C}{\mathbb C}
\newcommand{\Q}{\mathbb Q}
\newcommand{\bigO}{\mathcal O}
\newcommand{\smallo}{o}
\newcommand{\Acal}{\mathcal A}
\newcommand{\Pcal}{\mathcal P}
\newcommand{\Ncal}{\mathcal N}
\newcommand{\Rcal}{\mathcal R}
\newcommand{\Jcal}{\mathcal J}
\newcommand{\Bhat}{\widehat{B}}
\newcommand{\coeff}[2]{[#1^{#2}]}
\newcommand{\multi}[1]{\boldsymbol{#1}}
\newcommand{\abs}[1]{\lvert#1\rvert}
\newcommand{\norm}[1]{\lVert#1\rVert}
\newcommand{\floor}[1]{\left\lfloor#1\right\rfloor}
\newcommand{\ceil}[1]{\left\lceil#1\right\rceil}

\title{\bfseries Rademacher Transseries for the Partition Numbers\\[2mm]
       and Their Inverse\\[2mm]
       \large All-order arithmetic sectors, coefficient formulae,\\
       Lambert--Lagrange reversion, and phase-locked inversion}
\author{Prepared for the ProveIt/Fabius Function research corpus}
\date{3 September 2026}

\begin{document}
\maketitle

\begin{abstract}
Let $p(n)$ be the unrestricted partition number.  The exact Rademacher series is
reorganized here as a complete arithmetic transseries.  In the natural shifted
variable
\[
 x=n-\frac1{24},\qquad \mu=\pi\sqrt{\frac{2x}{3}},
\]
each fixed Rademacher sector is exactly an exponential multiplied by a
\emph{two-term} rational prefactor.  Re-expansion in the conventional variable
$n^{-1/2}$ produces, for every sector $(k,\sigma)$ with $k\geq1$ and
$\sigma\in\{+1,-1\}$, a convergent all-order coefficient series.  We give three
equivalent general formulae for its coefficients: an analytic generating
function, a complete-Bell-polynomial expression, and the finite sum
\[
 C_{k,\sigma;r}=\frac1{(-4\sqrt6)^r}
 \sum_{j=0}^{\lfloor(r+1)/2\rfloor}
 \binom{r+1}{j}\frac{r+1-j}{(r+1-2j)!}
 \left(\frac{\sigma\pi}{6k}\right)^{r-2j}.
\]
The case $(k,\sigma)=(1,+1)$ is the known O'Sullivan coefficient formula;
the displayed expression supplies its sector-wide extension.

Because $p$ is a sequence, ``the inverse'' first requires a convention.  We
construct a canonical real-analytic Rademacher interpolation $\Pcal(\nu)$ that
agrees with $p(n)$ at every integer and is strictly increasing for all
sufficiently large $\nu$.  Its inverse $\Ncal(y)$ therefore determines the
generalized inverse of the sequence by
$p_{\Z}^{\langle-1\rangle}(y)=\ceil{\Ncal(y)}$ for large $y$.
The dominant inverse is anchored by the exact Lambert carrier
\[
 \rho(y)=-2W_{-1}\!\left(-\frac12\sqrt{\frac{\kappa}{y}}\right),
 \qquad \kappa=\frac{\pi^2}{6\sqrt3},
\]
and the full algebraic correction is generated by Lagrange--B\"urmann and
ordinary Bell polynomials.  Its coefficients are given both by finite
coefficient extraction and by a triangular recurrence.  The complete
nonperturbative inverse is then obtained by a second Lagrange inversion.  A
finite multi-index formula gives every product-sector coefficient, including
all derivatives of the periodic Dedekind-sum amplitudes.

Two structural points are essential.  First, the positive Rademacher actions
$1-1/k$ accumulate at $1$, so the completed object is an arithmetically indexed,
paired-summation transseries rather than a finite-rank Hahn series.  Second, the
Lambert carrier and the dominant-core inverse differ by
$3/\pi^2+o(1)$ in the original index.  This constant must be resummed before the
periodic amplitudes are evaluated; otherwise their phases are wrong at leading
nonperturbative order.  The leading inverse correction is the parity chirp
\[
 \Ncal(y)-N_0(y)\sim
 -\frac{3\sqrt\kappa}{\pi^2\sqrt2}\,
 \frac{\cos(\pi N_0(y))}{\sqrt y}.
\]
Explicit remainder transfer principles, implementation algorithms, and
high-precision numerical checks are included.
\end{abstract}

\tableofcontents
\clearpage

\section{Introduction}

The unrestricted partition function $p(n)$ counts non-increasing finite
sequences of positive integers with sum $n$; it is OEIS A000041
\cite{OEISA000041}.  Its generating function is
Euler's product
\begin{equation}\label{eq:euler-product}
 \sum_{n\geq0}p(n)q^n=\prod_{m\geq1}\frac1{1-q^m},\qquad \abs q<1.
\end{equation}
Hardy and Ramanujan proved the celebrated first-order estimate
\cite{HardyRamanujan1918}
\begin{equation}\label{eq:HR}
 p(n)\sim \frac{1}{4\sqrt3\,n}
 \exp\!\left(\pi\sqrt{\frac{2n}{3}}\right).
\end{equation}
Rademacher's refinement is not merely an asymptotic series: it is a convergent
sum indexed by positive integers $k$, and each $k$ carries a finite exponential
sum built from Dedekind sums \cite{Rademacher1937PNAS,Rademacher1938,Rademacher1943,Lehmer1939,KongTeo2023}.  Consequently the genuinely complete asymptotic
object has three interacting structures:
\begin{enumerate}[label=(\roman*)]
\item an algebraic half-power expansion in $n^{-1/2}$;
\item exponentially small Rademacher sectors with actions $1\mp1/k$;
\item periodic arithmetic amplitudes, which become chirped functions of
      $\log y$ after inversion.
\end{enumerate}
Keeping only item (i) gives the usual Poincare expansion.  It suppresses all
arithmetic oscillation and therefore cannot describe the exponentially accurate
inverse.

This article develops the direct and inverse expansions in a form compatible
with the coefficient calculus based on Bell polynomials, Faa di Bruno
composition, and Lagrange--B\"urmann inversion
\cite{CombinatorialCoefficientCalculus,PolynomialLogTransseries}.  Formal identities are kept
separate from analytic error statements.  In particular, an exact convergent
Rademacher expression supplies the nonperturbative completion, while finite
coefficient extraction gives every formal coefficient.

\subsection{What is meant by ``full''}

For fixed $k$, the two exponentials $\e^{\pm\mu/k}$ are perfectly ordinary
transseries sectors.  Across all $k$, however, the actions
\[
 1-\frac12,\ 1-\frac13,\ 1-\frac14,\ldots
\]
approach $1$ from below, while the actions $1+1/k$ approach $1$ from above.
The latter set has no least element after restricting to an infinite tail.  In
addition, the two signs for a given large $k$ must be paired before summing: the
large-$k$ convergence is produced by cancellation between them.  Thus the
complete object is not a classical grid-based transseries with globally
well-ordered support.

We use the following precise replacement.

\begin{definition}[Rademacher transseries]\label{def:rademacher-transseries}
A Rademacher transseries is a compatible family of finite-$K$ resolved sector
expansions, together with the paired convergent Rademacher tail.  Every action
cutoff strictly below $1$ contains only finitely many resolved primary sectors.
At and beyond the accumulation action $1$, the paired tail is retained as an
analytic coefficient, or is resolved transasymptotically by increasing $K$.
\end{definition}

This definition preserves all information in the exact series and avoids an
illegitimate rearrangement of a conditionally structured infinite sector sum.

\subsection{Main results at a glance}

The direct expansion has an exceptionally simple shifted form.  Put
\begin{equation}\label{eq:basic-constants}
 c:=\pi\sqrt{\frac23}=\frac{\pi\sqrt6}{3},\qquad
 x:=n-\frac1{24},\qquad \mu:=c\sqrt x,
 \qquad \kappa:=\frac{c^2}{4\sqrt3}=\frac{\pi^2}{6\sqrt3}.
\end{equation}
For every $k\geq1$ and $\sigma=\pm1$, the $(k,\sigma)$ sector is
\begin{equation}\label{eq:sector-preview}
 P_{k,\sigma}(n)=
 \frac{A_k(n)}{4\sqrt{3k}\,x}\,
 \e^{\sigma\mu/k}\left(1-\frac{\sigma k}{\mu}\right).
\end{equation}
This is exact.  Re-expanding it about unshifted $n=\infty$ gives the all-order
coefficients stated in the abstract.

For inversion, the dominant positive sector is
\begin{equation}\label{eq:dominant-preview}
 D(\mu)=\kappa\frac{\e^\mu}{\mu^2}
 \left(1-\frac1\mu\right).
\end{equation}
The Lambert carrier solves the simpler equation
$y=\kappa\e^\rho/\rho^2$.  The exact dominant-core inverse $\mu_0$ solves
$y=D(\mu_0)$ and has a convergent expansion
\begin{equation}\label{eq:mu0-preview}
 \mu_0=\rho+\frac1\rho+\frac{5}{2\rho^2}
 +\frac{13}{3\rho^3}+\frac{53}{12\rho^4}-\frac{14}{5\rho^5}+\cdots.
\end{equation}
All remaining sectors are reverted by a universal second Lagrange formula.

\begin{resultbox}
The article gives explicit general coefficient formulae at every level:
\begin{enumerate}[label=\arabic*.]
\item $C_{k,\sigma;r}$ for every direct Rademacher sector, in finite-sum and
      Bell-polynomial form;
\item $P_m(\log L)$ for the elementary logarithmic expansion of the Lambert
      carrier, in signed-Stirling-number form;
\item $a_m$ for the dominant-core correction, by finite Lagrange extraction and
      a triangular ordinary-Bell recurrence;
\item $V_{\multi m}$ for every product of nonperturbative inverse sectors, by a
      finite multi-index Lagrange formula.
\end{enumerate}
\end{resultbox}

\section{Arithmetic and coefficient conventions}

\subsection{Dedekind sums and Rademacher amplitudes}

For real $u$, let the sawtooth function be
\begin{equation}\label{eq:sawtooth}
 ((u)):=
 \begin{cases}
 u-\floor{u}-\frac12,&u\notin\Z,\\
 0,&u\in\Z.
 \end{cases}
\end{equation}
For $k\geq1$ and $h\in\Z$, the Dedekind sum is
\begin{equation}\label{eq:dedekind}
 s(h,k):=\sum_{r=1}^{k-1}
 \left(\left(\frac rk\right)\right)
 \left(\left(\frac{hr}{k}\right)\right).
\end{equation}
The Rademacher amplitude at an integer $n$ may be written either as a complex
exponential sum or, because conjugate terms pair, as the real cosine sum
\begin{align}
 A_k(n)
 &=\sum_{\substack{0\leq h<k\\(h,k)=1}}
 \exp\!\left(\pi\ii s(h,k)-\frac{2\pi\ii nh}{k}\right),
 \label{eq:A-complex}\\
 &=\sum_{\substack{1\leq h\leq k\\(h,k)=1}}
 \cos\!\left(\pi s(h,k)-\frac{2\pi nh}{k}\right).
 \label{eq:A-cos}
\end{align}
The convention gives $A_1(n)=1$.  The elementary bound
\begin{equation}\label{eq:A-trivial-bound}
 \abs{A_k(n)}\leq \varphi(k)\leq k
\end{equation}
is sufficient for the convergence estimates used below.

The first amplitudes are
\begin{align}
 A_1(n)&=1,\nonumber\\
 A_2(n)&=(-1)^n,\nonumber\\
 A_3(n)&=2\cos\!\left(\frac{\pi}{18}-\frac{2\pi n}{3}\right),
 \label{eq:first-amplitudes}\\
 A_4(n)&=2\cos\!\left(\frac{\pi}{8}-\frac{\pi n}{2}\right).
 \nonumber
\end{align}
No asymptotic averaging of these amplitudes is performed: their arithmetic
phase is part of the answer.

\subsection{Bell polynomials}

The partial exponential Bell polynomials $B_{n,j}$ are defined by
\begin{equation}\label{eq:exp-bell-def}
 B_{n,j}(x_1,x_2,\ldots)=
 n!\!\sum_{\substack{m_1+m_2+\cdots=j\\
                     m_1+2m_2+\cdots=n}}
 \prod_{r\geq1}\frac1{m_r!}
 \left(\frac{x_r}{r!}\right)^{m_r}.
\end{equation}
The complete polynomial is $B_n=\sum_{j=0}^nB_{n,j}$.  The basic exponential
identity is
\begin{equation}\label{eq:bell-exp}
 \exp\!\left(\sum_{r\geq1}x_r\frac{t^r}{r!}\right)
 =\sum_{n\geq0}B_n(x_1,\ldots,x_n)\frac{t^n}{n!}.
\end{equation}

For an ordinary series $U(t)=\sum_{r\geq1}u_rt^r$, write
\begin{equation}\label{eq:ordinary-bell}
 \Bhat_{n,j}(u_1,u_2,\ldots):=[t^n]U(t)^j.
\end{equation}
Thus $\Bhat_{0,0}=1$, $\Bhat_{n,0}=0$ for $n>0$, and
\begin{equation}\label{eq:ordinary-exp-scaling}
 \Bhat_{n,j}(u)=\frac{j!}{n!}
 B_{n,j}(1!u_1,2!u_2,\ldots).
\end{equation}

\subsection{Signed Stirling numbers}

We use signed Stirling numbers of the first kind $s(m,j)$:
\begin{equation}\label{eq:stirling-def}
 z(z-1)\cdots(z-m+1)=\sum_{j=0}^m s(m,j)z^j.
\end{equation}
Equivalently,
\begin{equation}\label{eq:log-stirling}
 \log(1+z)^j=j!\sum_{m\geq j}s(m,j)\frac{z^m}{m!}.
\end{equation}

\subsection{Coefficient extraction and Lagrange--B\"urmann}

If $g(z)=z\phi(g(z))$ with $\phi(0)$ invertible, then
\begin{equation}\label{eq:LB-basic}
 [z^m]g(z)=\frac1m[w^{m-1}]\phi(w)^m,
 \qquad m\geq1.
\end{equation}
We repeatedly introduce a bookkeeping parameter $z$, apply
\eqref{eq:LB-basic}, and finally set $z=1$.  In each application the resulting
coefficient at a prescribed algebraic or exponential order is a finite sum.

\section{The exact Rademacher decomposition}

\subsection{Rademacher's convergent series}

\begin{theorem}[Rademacher \cite{Rademacher1937PNAS,Rademacher1938,EberlPaulson2026}]\label{thm:rademacher}
For every integer $n\geq1$,
\begin{equation}\label{eq:rademacher}
 p(n)=\frac1{\pi\sqrt2}\sum_{k=1}^{\infty}
 A_k(n)\sqrt{k}\,\frac{\dd}{\dd n}
 \left[
 \frac{\sinh\!\left(\frac{c}{k}\sqrt{n-\frac1{24}}\right)}
      {\sqrt{n-\frac1{24}}}
 \right],
\end{equation}
where $c=\pi\sqrt{2/3}$ and the series converges when summed in increasing
$k$.
\end{theorem}

The square root of $k$ in \eqref{eq:rademacher} is important.  Some modern
normalizations absorb it into the definition of the arithmetic sum; throughout
this article $A_k$ means \eqref{eq:A-complex}--\eqref{eq:A-cos}.

\subsection{Exact two-exponential sectors}

\begin{lemma}[Elementary derivative]\label{lem:elementary-derivative}
For $a>0$ and $x>0$,
\begin{equation}\label{eq:elementary-derivative}
 \frac{\dd}{\dd x}\left(x^{-1/2}\sinh(a\sqrt x)\right)
 =\frac1{4x^{3/2}}
 \left[(a\sqrt x-1)\e^{a\sqrt x}
      +(a\sqrt x+1)\e^{-a\sqrt x}\right].
\end{equation}
\end{lemma}

\begin{proof}
Differentiate the product $x^{-1/2}\sinh(a\sqrt x)$, replace hyperbolic
functions by exponentials, and collect the two signs.  No asymptotic expansion
is involved.
\end{proof}

\begin{theorem}[Exact sector formula]\label{thm:exact-sectors}
Let $x=n-1/24$ and $\mu=c\sqrt x$.  Define, for $\sigma=\pm1$,
\begin{equation}\label{eq:Pks}
 P_{k,\sigma}(n):=
 \frac{A_k(n)}{4\sqrt{3k}\,x}\,
 \e^{\sigma\mu/k}\left(1-\frac{\sigma k}{\mu}\right).
\end{equation}
Then Rademacher's series is exactly
\begin{equation}\label{eq:p-sector-sum}
 p(n)=\sum_{k\geq1}\bigl(P_{k,+}(n)+P_{k,-}(n)\bigr),
\end{equation}
where the two signs are paired before the sum over $k$ is taken.
Equivalently, with
\begin{equation}\label{eq:g-def}
 g(z):=\cosh z-\frac{\sinh z}{z},\qquad g(0):=0,
\end{equation}
we have the paired form
\begin{equation}\label{eq:paired-sector}
 P_{k,+}(n)+P_{k,-}(n)
 =\frac{A_k(n)}{2\sqrt{3k}\,x}\,g\!\left(\frac\mu k\right).
\end{equation}
\end{theorem}

\begin{proof}
Insert \eqref{eq:elementary-derivative} with $a=c/k$ into
\eqref{eq:rademacher}.  Since
\[
 \frac{\sqrt{k}}{\pi\sqrt2}\frac1{4x^{3/2}}
 \frac{\mu}{k}
 =\frac1{4\sqrt{3k}\,x},
\]
the two exponentials give \eqref{eq:Pks}.  Their sum is
\eqref{eq:paired-sector}.
\end{proof}

\begin{remark}[Why pairing is mandatory]\label{rem:pairing}
For fixed $k$, splitting into $P_{k,+}$ and $P_{k,-}$ is harmless.  For
$k\to\infty$, each split term contains a component proportional to $k/\mu$;
the cancellation of those components occurs only in their sum.  The paired
kernel satisfies $g(\mu/k)=\mu^2/(3k^2)+O(k^{-4})$, which restores convergence.
Thus \eqref{eq:p-sector-sum} specifies an order of summation, not an absolutely
rearrangeable double series.
\end{remark}

\subsection{Dominant sector}

The positive $k=1$ sector is
\begin{equation}\label{eq:dominant-sector}
 P_{1,+}(n)=\frac{\e^\mu}{4\sqrt3\,x}
 \left(1-\frac1\mu\right)
 =\kappa\frac{\e^\mu}{\mu^2}\left(1-\frac1\mu\right).
\end{equation}
The negative $k=1$ sector is smaller by $\e^{-2\mu}$ relative to it.  The
positive $k=2$ sector is smaller by $\e^{-\mu/2}$ and is the first arithmetic
correction.  This immediately explains why the shifted approximation
\eqref{eq:dominant-sector} is exponentially more accurate than the elementary
Hardy--Ramanujan formula \eqref{eq:HR}.

\section{The complete direct Rademacher transseries}

\subsection{Resolved actions and exact relative factorization}

Let
\begin{equation}\label{eq:N-of-mu}
 N(\mu):=\frac1{24}+\frac{3\mu^2}{2\pi^2},
\end{equation}
so that $\mu(N(\mu))=\mu$.  The same notation will later be used for the real
interpolation.  Relative to \eqref{eq:dominant-sector}, define
\begin{align}
 \alpha_{k,\sigma}&:=1-\frac{\sigma}{k},
 \label{eq:actions}\\
 r_{k,\sigma}(\mu)&:=
 \frac{A_k(N(\mu))}{\sqrt{k}}
 \frac{1-\sigma k/\mu}{1-1/\mu}.
 \label{eq:rks}
\end{align}
At integer arguments, $A_k(N(\mu))$ is interpreted on the corresponding
lattice point.  Then, for any finite $K$,
\begin{equation}\label{eq:finite-resolved}
 \frac{p(n)}{P_{1,+}(n)}
 =1+r_{1,-}(\mu)\e^{-2\mu}
 +\sum_{k=2}^{K}\sum_{\sigma=\pm1}
 r_{k,\sigma}(\mu)\e^{-\alpha_{k,\sigma}\mu}
 +\Rcal_{>K}(\mu).
\end{equation}
The actions are
\begin{equation}\label{eq:action-list}
 \alpha_{k,+}=1-\frac1k,\qquad
 \alpha_{k,-}=1+\frac1k.
\end{equation}

\subsection{The paired tail and action accumulation}

The exact non-dominant relative correction can be written without any split
large-$k$ sectors:
\begin{equation}\label{eq:R-exact}
 \Rcal(\mu)=
 \frac{1+1/\mu}{1-1/\mu}\e^{-2\mu}
 +\frac{2\e^{-\mu}}{1-1/\mu}
 \sum_{k=2}^{\infty}\frac{A_k(N(\mu))}{\sqrt{k}}
 g\!\left(\frac\mu k\right).
\end{equation}
Thus
\begin{equation}\label{eq:exact-factorization}
 p(n)=\kappa\frac{\e^\mu}{\mu^2}
 \left(1-\frac1\mu\right)\bigl(1+\Rcal(\mu)\bigr).
\end{equation}
Equation \eqref{eq:R-exact} is the analytic completion of the resolved
transseries.

\begin{lemma}[A useful kernel bound]\label{lem:g-bound}
For $z\geq0$,
\begin{equation}\label{eq:g-bound}
 0\leq g(z)\leq \frac{z^2}{3}\cosh z
 \leq \frac{z^2}{3}\e^z.
\end{equation}
\end{lemma}

\begin{proof}
The Taylor series is
\[
 g(z)=\sum_{m\geq1}\frac{2m}{(2m+1)!}z^{2m}.
\]
Since
$2m/(2m+1)!\leq1/(3(2m-2)!)$ for $m\geq1$, termwise comparison gives the
first upper bound.  The second is immediate.
\end{proof}

\begin{theorem}[Explicit finite-$K$ tail bound]\label{thm:tail-bound}
Let $\mu>1$ and $K\geq1$.  For the real cosine amplitudes, and hence at every
integer $n$, the paired tail in \eqref{eq:finite-resolved} satisfies
\begin{equation}\label{eq:tail-bound}
 \abs{\Rcal_{>K}(\mu)}
 \leq
 \frac{4\mu^2}{3(1-1/\mu)\sqrt K}
 \exp\!\left[-\left(1-\frac1{K+1}\right)\mu\right].
\end{equation}
\end{theorem}

\begin{proof}
By \eqref{eq:A-trivial-bound}, \eqref{eq:g-bound}, and
$\sum_{k>K}k^{-3/2}\leq2/\sqrt K$,
\begin{align*}
 \abs{\Rcal_{>K}}
 &\leq \frac{2\e^{-\mu}}{1-1/\mu}
 \sum_{k>K}\sqrt{k}\,g(\mu/k)\\
 &\leq \frac{2\mu^2\e^{-\mu}}{3(1-1/\mu)}
 \sum_{k>K}k^{-3/2}\e^{\mu/k}\\
 &\leq \frac{4\mu^2}{3(1-1/\mu)\sqrt K}
 \e^{-\mu+\mu/(K+1)}.
\end{align*}
\end{proof}

\begin{cautionbox}
The actions $1-1/k$ form an increasing sequence with limit $1$, but the actions
$1+1/k$ decrease to $1$.  Hence the complete primary action set is not globally
well ordered.  Formula \eqref{eq:tail-bound} is the correct replacement for a
formal assertion that an infinite list of split sectors may be freely sorted.
For every cutoff $\eta<1$, only finitely many positive sectors satisfy
$\alpha_{k,+}\leq\eta$; at the accumulation action, one uses the paired tail.
\end{cautionbox}

\section{All-order coefficients in the unshifted variable}

The shifted sector formula \eqref{eq:Pks} is exact and elementary.  It is often
convenient, however, to normalize by the conventional Hardy--Ramanujan
exponential $\e^{\sigma c\sqrt n/k}$ and expand in $n^{-1/2}$.  This section
gives all coefficients in three equivalent forms.

\subsection{Analytic generating function}

Set
\begin{equation}\label{eq:t-beta}
 t:=n^{-1/2},\qquad \delta:=\frac1{24},\qquad
 \beta_{k,\sigma}:=\frac{\sigma c}{k}.
\end{equation}

\begin{theorem}[Sector generating function]\label{thm:sector-generating}
For every fixed $(k,\sigma)$,
\begin{equation}\label{eq:sector-unshifted}
 P_{k,\sigma}(n)=
 \frac{A_k(n)}{4\sqrt{3k}\,n}
 \e^{\sigma c\sqrt n/k}
 \sum_{r\geq0}C_{k,\sigma;r}\,n^{-r/2},
\end{equation}
where
\begin{equation}\label{eq:C-generating}
 \boxed{
 \sum_{r\geq0}C_{k,\sigma;r}t^r=
 \frac{\displaystyle
 \exp\!\left[
 \frac{\sigma c}{kt}
 \left(\sqrt{1-\frac{t^2}{24}}-1\right)
 \right]}
 {\displaystyle 1-\frac{t^2}{24}}
 \left(1-
 \frac{\sigma kt}{c\sqrt{1-\frac{t^2}{24}}}
 \right).}
\end{equation}
The right-hand side is analytic at $t=0$ and has radius of analyticity at least
$\sqrt{24}$.  Consequently \eqref{eq:sector-unshifted} is convergent for every
integer $n\geq1$ when $k$ and $\sigma$ are fixed.
\end{theorem}

\begin{proof}
Write $x=n(1-t^2/24)$ and
$\mu=(c/t)\sqrt{1-t^2/24}$ in \eqref{eq:Pks}.  Factor out
$\e^{\sigma c/(kt)}$ and $1/n$.  The apparent singularity in the remaining
exponent is removable because
$\sqrt{1-t^2/24}-1=O(t^2)$.  The nearest branch points are at
$t=\pm\sqrt{24}$.
\end{proof}

\subsection{Finite closed form}

Define
\begin{equation}\label{eq:bks}
 b_{k,\sigma}:=\frac{\sigma\pi}{6k}.
\end{equation}

\begin{theorem}[Finite coefficient formula]\label{thm:C-finite}
For every $r\geq0$,
\begin{equation}\label{eq:C-finite}
 \boxed{
 C_{k,\sigma;r}=\frac1{(-4\sqrt6)^r}
 \sum_{j=0}^{\floor{(r+1)/2}}
 \binom{r+1}{j}\frac{r+1-j}{(r+1-2j)!}
 b_{k,\sigma}^{\,r-2j}.}
\end{equation}
\end{theorem}

\begin{proof}
A complete residue derivation is given in \cref{app:finite-sum-proof}.  The key
substitution is $u=-t/(4\sqrt6)$ followed by the Catalan change of variable
$T=u(1+T^2)$.  Lagrange residue transfer reduces the coefficient to
\[
 \frac{r+1}{b_{k,\sigma}}
 [T^{r+1}]\e^{b_{k,\sigma}T}(1+T^2)^r,
\]
whose binomial expansion is exactly \eqref{eq:C-finite}.
\end{proof}

The first coefficients, with $b=b_{k,\sigma}$, are
\begin{align}
 C_0&=1,\nonumber\\
 C_1&=-\frac{\sqrt6\,(b^2+2)}{24b},\nonumber\\
 C_2&=\frac{b^2+12}{192},\nonumber\\
 C_3&=-\frac{\sqrt6\,(b^4+36b^2+72)}{13824b},
 \label{eq:C-first}\\
 C_4&=\frac{b^4+80b^2+720}{221184},\nonumber\\
 C_5&=-\frac{\sqrt6\,(b^6+150b^4+3600b^2+7200)}{26542080b}.
 \nonumber
\end{align}

\subsection{Complete-Bell-polynomial form}

Let
\begin{equation}\label{eq:h-coeffs}
 h_{2m-1}^{(k,\sigma)}:=
 \frac{\sigma c}{k}\binom{1/2}{m}\left(-\frac1{24}\right)^m,
 \qquad h_{2m}^{(k,\sigma)}:=0.
\end{equation}
Define
\begin{equation}\label{eq:Dlambda}
 D_r^{(\lambda;k,\sigma)}:=
 \sum_{q=0}^{\floor{r/2}}
 \frac{(\lambda)_q}{q!\,24^q}\,
 \frac{B_{r-2q}
 \bigl(1!h_1^{(k,\sigma)},\ldots,(r-2q)!h_{r-2q}^{(k,\sigma)}\bigr)}
 {(r-2q)!},
\end{equation}
with $D_r^{(\lambda)}=0$ for $r<0$.

\begin{proposition}[Bell-polynomial coefficient calculus]\label{prop:C-bell}
The sector coefficients are
\begin{equation}\label{eq:C-bell}
 \boxed{
 C_{k,\sigma;r}=D_r^{(1;k,\sigma)}
 -\frac{\sigma k}{c}\,D_{r-1}^{(3/2;k,\sigma)}.}
\end{equation}
\end{proposition}

\begin{proof}
The exponent in \eqref{eq:C-generating} is
$\sum_{j\geq1}h_jt^j$.  Apply \eqref{eq:bell-exp}, multiply first by
$(1-t^2/24)^{-1}$ and then by
$1-(\sigma k/c)t(1-t^2/24)^{-1/2}$.  Cauchy convolution gives
\eqref{eq:Dlambda}--\eqref{eq:C-bell}.
\end{proof}

Formula \eqref{eq:C-bell} is often preferable in symbolic systems because it
uses only standard composition primitives.  Formula \eqref{eq:C-finite} is
usually faster for numerical coefficient generation.

\subsection{Recovery of the ordinary Poincare expansion}

For the dominant positive sector, $b_{1,+}=\pi/6$.  Set
\begin{equation}\label{eq:omega-def}
 \omega_r:=C_{1,+;r}.
\end{equation}
Then \eqref{eq:C-finite} becomes
\begin{equation}\label{eq:omega}
 \omega_r=\frac1{(-4\sqrt6)^r}
 \sum_{j=0}^{\floor{(r+1)/2}}
 \binom{r+1}{j}\frac{r+1-j}{(r+1-2j)!}
 \left(\frac\pi6\right)^{r-2j},
\end{equation}
which is the O'Sullivan formula \cite{OSullivan2023,BanerjeeEtAl2022}.  Since every other fixed-$k$ sector is
exponentially small relative to $P_{1,+}$, one obtains, for each fixed $R$,
\begin{equation}\label{eq:poincare}
 p(n)=\frac{\e^{c\sqrt n}}{4\sqrt3\,n}
 \left(\sum_{r=0}^{R-1}\omega_r n^{-r/2}
 +\bigO_R(n^{-R/2})\right).
\end{equation}

\begin{remark}[Convergent sector, asymptotic total]\label{rem:convergent-vs-asymptotic}
The infinite series $\sum_r\omega_r n^{-r/2}$ converges to the single sector
$P_{1,+}(n)$ for $n\geq1$; it does not converge to $p(n)$.  The difference
$p-P_{1,+}$ is beyond all algebraic orders and is supplied by the remaining
Rademacher sectors.  Thus \eqref{eq:poincare} is an asymptotic statement about
the complete partition function, even though its algebraic coefficient series
has a convergent sectorial sum.
\end{remark}

The first dominant coefficients are
\begin{align}
 \omega_0&=1,\nonumber\\
 \omega_1&=-\frac{\sqrt6(\pi^2+72)}{144\pi},\nonumber\\
 \omega_2&=\frac{\pi^2+432}{6912},\nonumber\\
 \omega_3&=-\frac{\sqrt6(\pi^4+1296\pi^2+93312)}{2985984\pi},
 \label{eq:omega-first}\\
 \omega_4&=\frac{\pi^4+2880\pi^2+933120}{286654464}.
 \nonumber
\end{align}

\section{A canonical inverse for a discrete sequence}

\subsection{Real-analytic Rademacher interpolation}

A sequence has no unique continuous inverse.  To retain the exponentially
small arithmetic information, an interpolation must also retain the
Rademacher amplitudes.  The cosine representation provides a natural choice.

\begin{definition}[Real Rademacher amplitude]\label{def:real-A}
For real or complex $\nu$, define
\begin{equation}\label{eq:A-real}
 \Acal_k(\nu):=
 \sum_{\substack{1\leq h\leq k\\(h,k)=1}}
 \cos\!\left(\pi s(h,k)-\frac{2\pi h\nu}{k}\right).
\end{equation}
Then $\Acal_k(n)=A_k(n)$ for every integer $n$.
\end{definition}
The first off-lattice amplitudes are
\begin{equation}\label{eq:real-low-amplitudes}
 \Acal_2(\nu)=\cos(\pi\nu),\qquad
 \Acal_3(\nu)=2\cos(\pi\nu)
 \cos\!\left(\frac\pi{18}+\frac{\pi\nu}{3}\right).
\end{equation}
At integer $\nu=n$, the second expression reduces to the simpler lattice form
in \eqref{eq:first-amplitudes}; that lattice-only simplification must not be used
when an inverse phase is evaluated at the generally nonintegral point $N_0(y)$.

For real $\nu>1/24$, set
\begin{equation}\label{eq:P-interpolation}
 \Pcal(\nu):=\frac1{\pi\sqrt2}\sum_{k\geq1}
 \Acal_k(\nu)\sqrt{k}\,\frac{\dd}{\dd \nu}
 \left[
 \frac{\sinh\!\left(\frac{c}{k}\sqrt{\nu-\frac1{24}}\right)}
      {\sqrt{\nu-\frac1{24}}}
 \right].
\end{equation}

\begin{proposition}[Interpolation and eventual monotonicity]\label{prop:interpolation}
The series \eqref{eq:P-interpolation} converges locally uniformly on
$\nu>1/24$ after the two signs are paired for each $k$.  It defines a
real-analytic function satisfying $\Pcal(n)=p(n)$ for every integer $n\geq1$.
Moreover, there exists $\nu_0$ such that $\Pcal'(\nu)>0$ for
$\nu\geq\nu_0$.
\end{proposition}

\begin{proof}
The paired $k$th term is \eqref{eq:paired-sector} with $A_k$ replaced by
$\Acal_k$.  On compact subsets, $\abs{\Acal_k}\leq k$ and
$g(\mu/k)=O(k^{-2})$, so the terms are $O(k^{-3/2})$ uniformly.  The same
argument applies after differentiation; derivatives of $\Acal_k$ have the same
$O(k)$ trivial bound.  This proves local uniform convergence and permits
termwise differentiation.  Equality at integers follows from
\cref{thm:rademacher}.

The dominant term is $D(\mu)$ from \eqref{eq:dominant-sector}.  Its logarithmic
derivative with respect to $\mu$ is
\begin{equation}\label{eq:F0prime-preview}
 1-\frac2\mu+\frac1{\mu(\mu-1)}=1+O(\mu^{-1}),
\end{equation}
and $\dd\mu/\dd\nu=c^2/(2\mu)>0$.  The differentiated non-dominant part is,
by the same estimates as \cref{thm:tail-bound}, a polynomial in $\mu$ times
$\e^{-\mu/2}$ relative to the dominant derivative.  Hence
$\Pcal'(\nu)=D'(\nu)(1+o(1))>0$.
\end{proof}

\subsection{Smooth and discrete inverses}

Let $\Ncal(y)$ denote the inverse branch of $\Pcal$ on
$[\nu_0,\infty)$.  For the sequence, define the generalized inverse
\begin{equation}\label{eq:integer-inverse}
 p_{\Z}^{\langle-1\rangle}(y):=
 \min\{n\in\N:p(n)\geq y\}.
\end{equation}

\begin{corollary}[Exact rounding relation]\label{cor:rounding}
For all sufficiently large real $y$,
\begin{equation}\label{eq:rounding}
 \boxed{
 p_{\Z}^{\langle-1\rangle}(y)=\ceil{\Ncal(y)}.}
\end{equation}
In particular, if $y=p(n)$ and $n$ is large, then $\Ncal(y)=n$.
\end{corollary}

\begin{proof}
On the interval of monotonicity, $\Pcal$ is increasing and agrees with $p$ at
integers.  If $p(n-1)<y\leq p(n)$, then $n-1<\Ncal(y)\leq n$.
\end{proof}

\begin{remark}[Dependence on interpolation]\label{rem:interpolation-dependence}
Other smooth interpolations can agree with $p(n)$ on the lattice and differ
between integers.  The algebraic inverse expansion is interpolation-independent;
nonperturbative phases are not.  Definition \ref{def:real-A} is canonical for
this article because it is obtained by extending the exact cosine amplitudes
already present in Rademacher's formula.  The integer generalized inverse is,
of course, independent of that choice after rounding.
\end{remark}

\section{Lambert carrier and complete algebraic inversion}

\subsection{Dominant-core equation}

For the interpolation, put
\begin{equation}\label{eq:Nmu-cont}
 N(\mu)=\frac1{24}+\frac{3\mu^2}{2\pi^2}
\end{equation}
and replace $A_k$ by $\Acal_k$ in \eqref{eq:R-exact}.  The exact inverse
equation is
\begin{equation}\label{eq:inverse-master}
 y=\kappa\frac{\e^\mu}{\mu^2}
 \left(1-\frac1\mu\right)(1+\Rcal(\mu)).
\end{equation}
Write
\begin{equation}\label{eq:L-def}
 L:=\log\frac{y}{\kappa}.
\end{equation}
The dominant-core phase function is
\begin{equation}\label{eq:F0}
 F_0(\mu):=\mu-2\log\mu+\log\left(1-\frac1\mu\right).
\end{equation}
It satisfies
\begin{equation}\label{eq:F0-derivatives}
 F_0'(\mu)=1-\frac2\mu+\frac1{\mu(\mu-1)},
 \qquad F_0'(\mu)=1+O(\mu^{-1}).
\end{equation}

\begin{definition}[Lambert carrier and dominant core]\label{def:carrier-core}
The Lambert carrier $\rho=\rho(y)$ is the large solution of
\begin{equation}\label{eq:rho-equation}
 \rho-2\log\rho=L.
\end{equation}
The dominant-core variable $\mu_0=\mu_0(y)$ is the large solution of
\begin{equation}\label{eq:mu0-equation}
 F_0(\mu_0)=L.
\end{equation}
Finally,
\begin{equation}\label{eq:N0-def}
 N_0(y):=\frac1{24}+\frac{3\mu_0(y)^2}{2\pi^2}.
\end{equation}
\end{definition}

\subsection{Exact Lambert representation}

\begin{proposition}[Lambert $W_{-1}$ carrier]\label{prop:rho-W}
For sufficiently large $y$,
\begin{equation}\label{eq:rho-W}
 \boxed{
 \rho(y)=-2W_{-1}\!\left(-\frac12
 \sqrt{\frac{\kappa}{y}}\right).}
\end{equation}
The branch $W_{-1}$ is forced by $\rho\to+\infty$ as $y\to+\infty$
\cite{CorlessEtAl1996,JeffreyEtAl1995}.
\end{proposition}

\begin{proof}
Equation \eqref{eq:rho-equation} is equivalent to
$\rho^2\e^{-\rho}=\kappa/y$.  Setting $w=-\rho/2$ gives
$w\e^w=-\tfrac12\sqrt{\kappa/y}$.  The principal branch tends to zero;
$W_{-1}$ tends to $-\infty$.
\end{proof}

\subsection{Elementary logarithmic expansion and signed-Stirling formula}

Let $\ell:=\log L$.  The carrier has the formal expansion
\begin{equation}\label{eq:rho-log-expansion}
 \rho=L+2\ell+\sum_{m\geq1}\frac{P_m(\ell)}{L^m}.
\end{equation}

\begin{theorem}[Closed formula for the Lambert polynomials]\label{thm:P-stirling}
For $m\geq1$,
\begin{equation}\label{eq:P-stirling}
 \boxed{
 P_m(\ell)=\frac1{m!}\sum_{j=1}^{m}
 2^j(j-1)!\,s(m,j)\binom{m}{j-1}
 (2\ell)^{m-j+1}.}
\end{equation}
The first four are
\begin{align}
 P_1(\ell)&=4\ell,\nonumber\\
 P_2(\ell)&=-4\ell(\ell-2),\nonumber\\
 P_3(\ell)&=\frac{8\ell}{3}(2\ell^2-9\ell+6),
 \label{eq:P-first}\\
 P_4(\ell)&=-\frac{8\ell}{3}(3\ell^3-22\ell^2+36\ell-12).
 \nonumber
\end{align}
\end{theorem}

\begin{proof}
Put $v=\rho-L-2\ell$ and $q=L^{-1}$.  Equation
\eqref{eq:rho-equation} becomes
\begin{equation}\label{eq:v-rho-fixed}
 v=2\log\bigl(1+q(2\ell+v)\bigr).
\end{equation}
Introduce $v=z\,2\log(1+q(2\ell+v))$.  Lagrange--B\"urmann gives
\[
 v=\sum_{j\geq1}\frac{2^j}{j}[w^{j-1}]
 \log(1+q(2\ell+w))^j.
\]
Using \eqref{eq:log-stirling}, the coefficient of $q^m$ is
\[
 \frac1{m!}\sum_{j=1}^m2^j(j-1)!s(m,j)
 [w^{j-1}](2\ell+w)^m,
\]
which is \eqref{eq:P-stirling}.
\end{proof}

\subsection{The dominant-core correction}

Write
\begin{equation}\label{eq:mu0-U}
 \mu_0=\rho+U(\rho^{-1}),\qquad
 U(s)=\sum_{m\geq1}a_ms^m.
\end{equation}

\begin{theorem}[All algebraic inverse coefficients]\label{thm:a-coefficients}
The series $U(s)$ is the unique analytic solution near $s=0$ of
\begin{equation}\label{eq:U-fixed}
 \boxed{
 U=3\log(1+sU)-\log(1-s+sU).}
\end{equation}
Its coefficients have the finite Lagrange representation
\begin{equation}\label{eq:a-lagrange}
 \boxed{
 a_m=[s^m]\sum_{q=1}^{m}\frac1q[w^{q-1}]
 \left(3\log(1+sw)-\log(1-s+sw)\right)^q.}
\end{equation}
Equivalently, with $\Bhat$ defined by \eqref{eq:ordinary-bell}, they satisfy
the triangular recurrence
\begin{align}
 a_m={}&3\sum_{q=1}^{\floor{m/2}}
 \frac{(-1)^{q+1}}{q}\,
 \Bhat_{m-q,q}(a_1,a_2,\ldots)
 \label{eq:a-recurrence}\\
 &+\sum_{q=1}^{m}\frac1q
 \sum_{j=0}^{\min(q,m-q)}(-1)^j\binom qj
 \Bhat_{m-q,j}(a_1,a_2,\ldots).
 \nonumber
\end{align}
\end{theorem}

\begin{proof}
Set $\mu_0=\rho+u$ in $F_0(\mu_0)=\rho-2\log\rho$.  With $s=1/\rho$,
\[
 0=u-2\log(1+su)+\log\left(1-\frac{s}{1+su}\right).
\]
Separating the last logarithm gives \eqref{eq:U-fixed}.  The analytic implicit
function theorem applies because the derivative with respect to $U$ at
$(s,U)=(0,0)$ is $1$.

Introduce a bookkeeping parameter $z$ on the right of \eqref{eq:U-fixed} and
apply \eqref{eq:LB-basic}; this proves \eqref{eq:a-lagrange}.  Finally expand
\[
 3\log(1+sU)=3\sum_{q\geq1}\frac{(-1)^{q+1}}q(sU)^q
\]
and
\[
 -\log(1-s+sU)=\sum_{q\geq1}\frac{s^q}{q}(1-U)^q.
\]
Extracting $[s^m]$ and using ordinary Bell polynomials gives
\eqref{eq:a-recurrence}.  Every term on the right involves only
$a_1,\ldots,a_{m-1}$.
\end{proof}

The first coefficients are
\begin{equation}\label{eq:a-table-inline}
\begin{split}
 a_1&=1,\quad a_2=\frac52,\quad a_3=\frac{13}{3},\quad
 a_4=\frac{53}{12},\quad a_5=-\frac{14}{5},\\
 a_6&=-\frac{763}{30},\quad a_7=-\frac{2179}{35},\quad
 a_8=-\frac{42289}{630},\quad a_9=\frac{132973}{1260},\quad
 a_{10}=\frac{3449711}{5040}.
\end{split}
\end{equation}
In particular,
\begin{equation}\label{eq:mu0-explicit}
 \mu_0=\rho+\rho^{-1}+\frac52\rho^{-2}
 +\frac{13}{3}\rho^{-3}+\frac{53}{12}\rho^{-4}
 -\frac{14}{5}\rho^{-5}+\cdots.
\end{equation}

\subsection{The algebraic inverse in the original index}

Substituting \eqref{eq:mu0-U} into \eqref{eq:N0-def} gives
\begin{align}
 N_0(y)
 ={}&\frac{3\rho^2}{2\pi^2}+\frac1{24}+\frac3{\pi^2}
 +\frac{15}{2\pi^2\rho}
 +\frac{29}{2\pi^2\rho^2}
 +\frac{83}{4\pi^2\rho^3}
 \label{eq:N0-rho}\\
 &+\frac{559}{40\pi^2\rho^4}
 -\frac{611}{20\pi^2\rho^5}
 -\frac{112459}{840\pi^2\rho^6}+\cdots.
 \nonumber
\end{align}
More generally, for $m\geq1$, the coefficient of $\rho^{-m}$ is
\begin{equation}\label{eq:N0-general-coeff}
 \frac{3}{2\pi^2}
 \left(2a_{m+1}+\sum_{j=1}^{m-1}a_ja_{m-j}\right).
\end{equation}

In elementary logarithms, retaining the first complete block gives
\begin{equation}\label{eq:N0-L}
\begin{split}
 N_0(y)={}&\frac{3}{2\pi^2}L^2
 +\frac{6}{\pi^2}L\log L
 +\frac{6}{\pi^2}(\log L)^2
 +\frac{12}{\pi^2}\log L\\
 &+\frac3{\pi^2}+\frac1{24}
 +\bigO\!\left(\frac{(\log L)^2}{L}\right),
 \qquad L=\log(y/\kappa).
\end{split}
\end{equation}
The exact Lambert form \eqref{eq:rho-W} plus the convergent $U$ series is
numerically preferable to expanding everything in $L$.

\subsection{Phase locking}

Let the uncorrected carrier index be
\begin{equation}\label{eq:Nrho}
 N_\rho(y):=\frac1{24}+\frac{3\rho(y)^2}{2\pi^2}.
\end{equation}
Then \eqref{eq:N0-rho} implies
\begin{equation}\label{eq:phase-shift}
 \boxed{
 N_0(y)-N_\rho(y)=\frac3{\pi^2}+O(\rho^{-1}).}
\end{equation}
The limiting shift is $3/\pi^2=0.3039635509\ldots$.

\begin{warning}[Do not expand the arithmetic phase about $N_\rho$]
The amplitudes $\Acal_k(N)$ are periodic on the $N$ scale.  Although
$\mu_0-\rho=O(\rho^{-1})$, the induced index change is
$N_0-N_\rho=O(1)$ because $\dd N/\dd\mu\asymp\rho$.  Taylor expansion of
$\Acal_k(N(\mu_0))$ around $\rho$ therefore does not begin with a small phase
perturbation.  The full algebraic core $N_0$ must be used as the base phase.
Only the subsequent exponentially small displacement can be expanded.
\end{warning}

\section{The full inverse transseries}

\subsection{One-stage universal Lagrange formula}

Let $\Rcal(\mu)$ be the exact paired correction \eqref{eq:R-exact}, with
$A_k$ replaced by the real amplitudes $\Acal_k$.  Put
\begin{equation}\label{eq:mu-rho-u}
 \mu=\rho+u.
\end{equation}
Taking logarithms in \eqref{eq:inverse-master} and subtracting
\eqref{eq:rho-equation} gives
\begin{equation}\label{eq:u-master}
 0=u-2\log\left(1+\frac u\rho\right)
 +\log\left(1-\frac1{\rho+u}\right)
 +\log\bigl(1+\Rcal(\rho+u)\bigr).
\end{equation}
Define
\begin{equation}\label{eq:Phi-rho}
 \Phi_\rho(w):=
 2\log\left(1+\frac w\rho\right)
 -\log\left(1-\frac1{\rho+w}\right)
 -\log\bigl(1+\Rcal(\rho+w)\bigr).
\end{equation}

\begin{theorem}[Complete one-stage inverse formula]\label{thm:one-stage-inverse}
As a formal Rademacher transseries, and analytically for sufficiently large
$y$ under the paired-summation convention,
\begin{equation}\label{eq:one-stage-inverse}
 \boxed{
 u=\sum_{q\geq1}\frac1q[w^{q-1}]\Phi_\rho(w)^q,\qquad
 \Ncal(y)=\frac1{24}+\frac{3(\rho+u)^2}{2\pi^2}.}
\end{equation}
Every prescribed coefficient of a finite-$K$, finite-action truncation receives
only finitely many contributions.
\end{theorem}

\begin{proof}
Introduce $u=z\Phi_\rho(u)$.  Lagrange--B\"urmann gives
$[z^q]u=q^{-1}[w^{q-1}]\Phi_\rho(w)^q$.  Setting $z=1$ gives the formal
identity.  For large $\rho$, the algebraic part is $O(\rho^{-1})$ and
$\Rcal=O(\rho^2\e^{-\rho/2})$; on a fixed small disk the map is a contraction,
which justifies the analytic branch.
\end{proof}

Formula \eqref{eq:one-stage-inverse} is compact but mixes algebraic and
exponential corrections.  The two-stage form below exposes every
nonperturbative coefficient explicitly and automatically respects phase
locking.

\subsection{Two-stage expansion about the dominant core}

Write
\begin{equation}\label{eq:mu-mu0-v}
 \mu=\mu_0+v.
\end{equation}
Set
\begin{equation}\label{eq:G-core}
 G(w):=
 \begin{cases}
 \displaystyle\frac{F_0(\mu_0+w)-F_0(\mu_0)}{w},&w\neq0,\\[3mm]
 F_0'(\mu_0),&w=0.
 \end{cases}
\end{equation}
Then the exact correction satisfies
\begin{equation}\label{eq:v-fixed-exact}
 v=-\frac{\log(1+\Rcal(\mu_0+v))}{G(v)}.
\end{equation}
Consequently,
\begin{equation}\label{eq:v-LB-exact}
 \boxed{
 v=\sum_{q\geq1}\frac1q[w^{q-1}]
 \left(-\frac{\log(1+\Rcal(\mu_0+w))}{G(w)}\right)^q.}
\end{equation}
This is the phase-locked complete inverse transseries: all arithmetic amplitudes
are based at $N_0=N(\mu_0)$.

\subsection{Resolved \texorpdfstring{finite-$K$}{finite-K} variables}

For a sector cutoff $K\geq2$, let
\begin{equation}\label{eq:JK}
 \Jcal_K:=\{(1,-)\}\cup
 \{(k,\sigma):2\leq k\leq K,\ \sigma=\pm1\}.
\end{equation}
For $j=(k,\sigma)\in\Jcal_K$, set
\begin{equation}\label{eq:E-B}
 \alpha_j=1-\frac\sigma k,
 \qquad E_j:=\e^{-\alpha_j\mu_0},
 \qquad B_j(w):=r_{k,\sigma}(\mu_0+w)\e^{-\alpha_jw},
\end{equation}
where now
\begin{equation}\label{eq:r-cont}
 r_{k,\sigma}(\mu)=
 \frac{\Acal_k(N(\mu))}{\sqrt{k}}
 \frac{1-\sigma k/\mu}{1-1/\mu}.
\end{equation}
Introduce independent formal variables $\multi E=(E_j)_{j\in\Jcal_K}$ and
\begin{equation}\label{eq:S-Q}
 S(w;\multi E):=\sum_{j\in\Jcal_K}E_jB_j(w),
 \qquad Q(w;\multi E):=\log(1+S(w;\multi E)).
\end{equation}
For a finite multi-index $\multi m=(m_j)$, write
$\abs{\multi m}=\sum_jm_j$, $\multi m!=\prod_jm_j!$, and
$\multi E^{\multi m}=\prod_jE_j^{m_j}$.

\begin{lemma}[Coefficients of the logarithm]\label{lem:qmulti}
For $\multi m\neq\multi0$,
\begin{equation}\label{eq:qmulti}
 [\multi E^{\multi m}]Q(w;\multi E)=
 (-1)^{\abs{\multi m}+1}
 \frac{(\abs{\multi m}-1)!}{\multi m!}
 \prod_{j\in\Jcal_K}B_j(w)^{m_j}.
\end{equation}
\end{lemma}

\begin{proof}
Only the term $S^{\abs{\multi m}}$ in
$\log(1+S)=\sum_{d\geq1}(-1)^{d+1}S^d/d$ contributes.  Its multinomial
coefficient is $\abs{\multi m}!/\multi m!$.
\end{proof}

\subsection{General multi-sector coefficient formula}

Expand
\begin{equation}\label{eq:v-multi}
 v=\sum_{\multi m\neq\multi0}V_{\multi m}(\mu_0)\,
 \multi E^{\multi m}.
\end{equation}

\begin{theorem}[All inverse-sector coefficients]\label{thm:Vmulti}
For every nonzero multi-index $\multi m$,
\begin{equation}\label{eq:Vmulti-compact}
 \boxed{
 V_{\multi m}=
 [\multi E^{\multi m}]
 \sum_{q=1}^{\abs{\multi m}}\frac1q[w^{q-1}]
 \left(-\frac{Q(w;\multi E)}{G(w)}\right)^q.}
\end{equation}
Equivalently, if
$q_{\multi a}(w):=[\multi E^{\multi a}]Q(w;\multi E)$ is given by
\eqref{eq:qmulti}, then
\begin{equation}\label{eq:Vmulti-expanded}
 V_{\multi m}=
 \sum_{q=1}^{\abs{\multi m}}\frac{(-1)^q}{q}
 [w^{q-1}]\frac1{G(w)^q}
 \sum_{\substack{\multi m_1+\cdots+\multi m_q=\multi m\\
                  \multi m_i\neq\multi0}}
 \prod_{i=1}^q q_{\multi m_i}(w).
\end{equation}
Both sums are finite.
\end{theorem}

\begin{proof}
Apply Lagrange--B\"urmann to
$v=z[-Q(v;\multi E)/G(v)]$ and extract $\multi E^{\multi m}$.  Since every
monomial of $Q$ has positive total $\multi E$-degree, only
$q\leq\abs{\multi m}$ contributes.  Expanding $Q^q$ gives
\eqref{eq:Vmulti-expanded}.
\end{proof}

This is the promised general coefficient formula for the full inverse
transseries.  It consists solely of finite coefficient extractions,
multinomial products, and derivatives at $w=0$; those derivatives may in turn
be written with complete Bell polynomials.

\subsection{Triangular recursion and low product orders}

Let $D_0:=F_0'(\mu_0)$.  Substitution of \eqref{eq:v-multi} into
\begin{equation}\label{eq:v-equation-expanded}
 D_0v+\sum_{r\geq2}\frac{F_0^{(r)}(\mu_0)}{r!}v^r
 +Q(v;\multi E)=0
\end{equation}
provides a triangular alternative to \cref{thm:Vmulti}:
\begin{equation}\label{eq:V-recursion}
 \boxed{
 V_{\multi m}=-\frac1{D_0}[\multi E^{\multi m}]
 \left[
 \sum_{r\geq2}\frac{F_0^{(r)}(\mu_0)}{r!}v^r
 +Q(v;\multi E)
 \right]_{\text{lower total degree}}.}
\end{equation}
The bracket at total degree $\abs{\multi m}$ contains no unknown
$V_{\multi m}$.

For a unit multi-index $\multi e_j$,
\begin{equation}\label{eq:V1}
 \boxed{
 V_{\multi e_j}=-\frac{r_j(\mu_0)}{F_0'(\mu_0)}.}
\end{equation}
Writing $B_j=B_j(0)$ and $B_j'=B_j'(0)$, the repeated-sector coefficient is
\begin{equation}\label{eq:V2same}
 V_{2\multi e_j}=-\frac1{D_0}
 \left[
 \frac12F_0''(\mu_0)V_{\multi e_j}^2
 +B_j'V_{\multi e_j}-\frac12B_j^2
 \right].
\end{equation}
For distinct $j$ and $\ell$,
\begin{equation}\label{eq:V2cross}
 V_{\multi e_j+\multi e_\ell}=-\frac1{D_0}
 \left[
 F_0''(\mu_0)V_{\multi e_j}V_{\multi e_\ell}
 +B_j'V_{\multi e_\ell}+B_\ell'V_{\multi e_j}-B_jB_\ell
 \right].
\end{equation}
These formulae display explicitly how nonlinear reversion mixes arithmetic
sectors.

\subsection{Coefficients for the inverse index}

Since
\begin{equation}\label{eq:N-v}
 \Ncal(y)=N_0(y)+\frac{3\mu_0}{\pi^2}v
 +\frac{3}{2\pi^2}v^2,
\end{equation}
the coefficient of $\multi E^{\multi m}$ in the inverse index is
\begin{equation}\label{eq:Nmulti}
 \boxed{
 \mathfrak N_{\multi m}=
 \frac{3\mu_0}{\pi^2}V_{\multi m}
 +\frac{3}{2\pi^2}
 \sum_{\substack{\multi a+\multi b=\multi m\\
                  \multi a,\multi b\neq\multi0}}
 V_{\multi a}V_{\multi b}.}
\end{equation}
Equations \eqref{eq:Vmulti-compact} and \eqref{eq:Nmulti} are an all-order,
closed coefficient calculus for the inverse.

\subsection{Primary inverse sectors}

At first nonperturbative order, \eqref{eq:V1} and \eqref{eq:N-v} give
\begin{equation}\label{eq:Delta-primary}
 \boxed{
 \Delta N_{k,\sigma}^{(1)}(y)=
 -\frac{3\mu_0}{\pi^2F_0'(\mu_0)}
 \frac{\Acal_k(N_0)}{\sqrt{k}}
 \frac{1-\sigma k/\mu_0}{1-1/\mu_0}
 \e^{-(1-\sigma/k)\mu_0}.}
\end{equation}
The first two positive sectors are
\begin{align}
 \Delta N_{2,+}^{(1)}
 &=-\frac{3\mu_0}{\pi^2F_0'(\mu_0)}
 \frac{\cos(\pi N_0)}{\sqrt2}
 \frac{1-2/\mu_0}{1-1/\mu_0}\e^{-\mu_0/2},
 \label{eq:Delta2}\\
 \Delta N_{3,+}^{(1)}
 &=-\frac{3\mu_0}{\pi^2F_0'(\mu_0)}
 \frac{2}{\sqrt3}\cos(\pi N_0)
 \cos\!\left(\frac\pi{18}+\frac{\pi N_0}{3}\right)
 \frac{1-3/\mu_0}{1-1/\mu_0}\e^{-2\mu_0/3}.
 \label{eq:Delta3}
\end{align}
Thus, before the accumulation action $1$,
\begin{equation}\label{eq:first-two-inverse}
 \Ncal(y)=N_0(y)+\Delta N_{2,+}^{(1)}(y)
 +\Delta N_{3,+}^{(1)}(y)
 +\bigO\!\left(\mu_0^3\e^{-3\mu_0/4}\right).
\end{equation}
The displayed remainder is a convenient coarse bound obtained from
\cref{thm:tail-bound}; individual resolved terms are generally smaller by
powers of $\mu_0$.

\subsection{Fractional powers of the inverse argument}

The dominant-core equation gives the exact conversion
\begin{equation}\label{eq:exp-to-y}
 \e^{-\alpha\mu_0}=
 \left(\frac\kappa y\right)^\alpha
 \mu_0^{-2\alpha}
 \left(1-\frac1{\mu_0}\right)^\alpha.
\end{equation}
Hence a primary positive $k$-sector has the scale
\begin{equation}\label{eq:positive-y-scale}
 \Delta N_{k,+}^{(1)}
 =y^{-(1-1/k)}\mu_0^{-1+2/k}
 \times\{\text{bounded arithmetic amplitude and a }\mu_0^{-1}\text{ series}\}.
\end{equation}
The staircase begins with
\[
 y^{-1/2},\qquad y^{-2/3}(\log y)^{-1/3},\qquad
 y^{-3/4}(\log y)^{-1/2},\ldots,
\]
and accumulates at the $y^{-1}$ level.

For $k=2$, \eqref{eq:Delta2} becomes
\begin{equation}\label{eq:Delta2-y-exact}
 \Delta N_{2,+}^{(1)}=
 -\frac{3\sqrt\kappa}{\pi^2\sqrt2}
 \frac{\cos(\pi N_0)}{\sqrt y}\,
 \frac{1-2/\mu_0}
 {F_0'(\mu_0)\sqrt{1-1/\mu_0}}.
\end{equation}
Therefore
\begin{equation}\label{eq:leading-inverse-chirp}
 \boxed{
 \Ncal(y)=N_0(y)-
 \frac{3\sqrt\kappa}{\pi^2\sqrt2}
 \frac{\cos(\pi N_0(y))}{\sqrt y}
 \left(1+O\!\left(\frac1{\log y}\right)\right)
 +\widetilde O(y^{-2/3}).}
\end{equation}
Here $\widetilde O$ permits powers of $\log y$.  The constant is
\begin{equation}\label{eq:leading-constant}
 \frac{3\sqrt\kappa}{\pi^2\sqrt2}
 =0.2094596846361762626\ldots.
\end{equation}
Because $N_0(y)\asymp(\log y)^2$, the cosine in
\eqref{eq:leading-inverse-chirp} is a quadratic logarithmic chirp, not a
log-periodic oscillation.

\subsection{Completion at the accumulation action}

For any $K$, \cref{thm:Vmulti} resolves all sectors in $\Jcal_K$ and all of
their products.  The omitted primary tail is kept in paired form:
\begin{equation}\label{eq:inverse-tail-completion}
 \Rcal_{>K}(\mu)=
 \frac{2\e^{-\mu}}{1-1/\mu}
 \sum_{k>K}\frac{\Acal_k(N(\mu))}{\sqrt{k}}
 g\!\left(\frac\mu k\right).
\end{equation}
One may insert \eqref{eq:inverse-tail-completion} as an additional analytic
symbol in \eqref{eq:v-LB-exact}.  Increasing $K$ produces a compatible
transasymptotic resolution, while \eqref{eq:tail-bound} controls the unresolved
part.  This is the precise meaning of the full inverse at and beyond action
$1$.

\section{Residual transfer and rigorous truncation control}

Formal coefficient matching does not by itself bound the error of an actual
inverse.  The needed analytic step is elementary but should be stated.

\begin{theorem}[Residual-to-inverse error transfer]\label{thm:residual-transfer}
Let
\begin{equation}\label{eq:H-full}
 H(\mu;y):=F_0(\mu)+\log(1+\Rcal(\mu))-\log(y/\kappa).
\end{equation}
Suppose $H(\mu;y)=0$ has a solution $\mu_*$ in an interval $I$, and
\begin{equation}\label{eq:derivative-lower}
 \inf_{\mu\in I}\abs{\partial_\mu H(\mu;y)}\geq m>0.
\end{equation}
If an approximation $\widetilde\mu\in I$ has residual
$\varepsilon=H(\widetilde\mu;y)$ and the line segment joining
$\widetilde\mu$ to $\mu_*$ lies in $I$, then
\begin{equation}\label{eq:mu-error}
 \abs{\widetilde\mu-\mu_*}\leq\frac{\abs\varepsilon}{m}.
\end{equation}
Consequently,
\begin{equation}\label{eq:N-error-transfer}
 \abs{N(\widetilde\mu)-N(\mu_*)}
 \leq \frac{3}{2\pi^2}
 \bigl(2\abs{\mu_*}+\abs{\widetilde\mu-\mu_*}\bigr)
 \frac{\abs\varepsilon}{m}.
\end{equation}
\end{theorem}

\begin{proof}
Apply the mean-value theorem to $H$ and then to
$N(\mu)=1/24+3\mu^2/(2\pi^2)$.
\end{proof}

For large $\mu$, $F_0'(\mu)=1+O(\mu^{-1})$ and
$\Rcal'(\mu)$ is exponentially small up to a polynomial factor, so one may
take, for example, $m=1/2$ beyond an effectively checkable threshold.  This
turns any formal truncation into a certificate once its residual is bounded.

\subsection{Algebraic truncation}

If
\begin{equation}\label{eq:U-M}
 U_M(\rho^{-1})=\sum_{m=1}^{M}a_m\rho^{-m},
\end{equation}
then the implicit analytic equation \eqref{eq:U-fixed} gives
\begin{equation}\label{eq:U-M-error}
 \mu_0-(\rho+U_M)=O(\rho^{-M-1}),
\end{equation}
and therefore
\begin{equation}\label{eq:N0-M-error}
 N_0-\left[\frac1{24}+\frac{3}{2\pi^2}(\rho+U_M)^2\right]
 =O(\rho^{-M}).
\end{equation}
To evaluate an arithmetic amplitude to $o(1)$ phase error, at least the
constant index shift in \eqref{eq:phase-shift} must be included.

\subsection{Sector truncation}

For a fixed $K$, substitute the resolved finite sum and paired tail bound into
\eqref{eq:H-full}.  Since $\log(1+z)=z+O(z^2)$ for small $z$, the direct bound
\eqref{eq:tail-bound} transfers to the inverse with one additional factor of
order $\mu$ in the original index.  In particular, for fixed $K$,
\begin{equation}\label{eq:inverse-K-bound}
 \Ncal(y)-\Ncal_K(y)
 =O_K\!\left(\mu_0^3
 \exp\!\left[-\left(1-\frac1{K+1}\right)\mu_0\right]\right),
\end{equation}
where $\Ncal_K$ is the inverse obtained from sectors through $k=K$ and the
$k=1$ negative sector.  Sharper bounds follow by using nontrivial estimates for
$A_k$ or by isolating the first omitted $k$.

\subsection{Certified recovery of an integer index}

For $y=p(n)$, any approximation $\widetilde N(y)$ satisfying
\begin{equation}\label{eq:round-cert}
 \abs{\widetilde N(y)-n}<\frac12
\end{equation}
recovers $n$ by nearest-integer rounding.  The dominant-core inverse $N_0(y)$
already has exponentially small error, so \eqref{eq:round-cert} holds for all
sufficiently large $n$.  Adding the first few primary sectors makes the margin
enormous even at moderate $n$.

\section{Algorithms}

All formulae in this article are constructive.  This section records concise
implementations independent of any one computer algebra system.

\subsection{Direct sector coefficients}

\begin{algorithm}[Finite-sum generation of $C_{k,\sigma;r}$]
Given $k\geq1$, $\sigma\in\{\pm1\}$, and $R$:
\begin{enumerate}[label=\arabic*.]
\item Set $b=\sigma\pi/(6k)$.
\item For $0\leq r\leq R$, compute
\[
 C_r=(-4\sqrt6)^{-r}\sum_{j=0}^{\floor{(r+1)/2}}
 \binom{r+1}{j}\frac{r+1-j}{(r+1-2j)!}b^{r-2j}.
\]
\item Use \eqref{eq:sector-unshifted}, pairing the two signs before summing over
      large $k$.
\end{enumerate}
The arithmetic complexity for one coefficient is $O(r)$ operations in the
coefficient field.
\end{algorithm}

\subsection{Lambert carrier polynomials}

\begin{algorithm}[Generation of $P_m(\ell)$]
Precompute signed Stirling numbers $s(m,j)$ by
\[
 s(m+1,j)=s(m,j-1)-m\,s(m,j),
 \qquad s(0,0)=1.
\]
Then apply \eqref{eq:P-stirling}.  This gives the full elementary-logarithmic
expansion \eqref{eq:rho-log-expansion}; for numerical work, evaluate the exact
$W_{-1}$ expression instead.
\end{algorithm}

\subsection{Dominant-core coefficients}

\begin{algorithm}[Triangular generation of $a_m$]
Initialize $a_1=1$.  For $m=2,3,\ldots,M$, compute the ordinary Bell
coefficients $\Bhat_{n,j}$ from the current truncated series
$U(s)=\sum_{r<m}a_rs^r$, and use \eqref{eq:a-recurrence}.  No nonlinear solve is
required at any stage.
\end{algorithm}

A direct symbolic alternative is the following pseudo-code.
\begin{lstlisting}
U = 0
for m = 1..M:
    Phi(w,s) = 3*log(1+s*w) - log(1-s+s*w)
    a[m] = coeff(s,m, sum(q=1..m, coeff(w,q-1,Phi(w,s)^q)/q))
    U = U + a[m]*s^m
\end{lstlisting}

\subsection{Full multi-sector inverse}

\begin{algorithm}[Coefficient of a prescribed multi-index]
Fix $K$ and a nonzero multi-index $\multi m$.
\begin{enumerate}[label=\arabic*.]
\item Compute $\rho$ from \eqref{eq:rho-W}, then $\mu_0$ from the convergent
      $U$ series or one Newton step on \eqref{eq:mu0-equation}.
\item Form $G(w)$, $B_j(w)$, and $q_{\multi a}(w)$ using
      \eqref{eq:G-core}, \eqref{eq:E-B}, and \eqref{eq:qmulti}.
\item Evaluate the finite sum \eqref{eq:Vmulti-expanded}; only
      $1\leq q\leq\abs{\multi m}$ and Taylor order $q-1$ are needed.
\item Convert $V_{\multi m}$ to the inverse-index coefficient with
      \eqref{eq:Nmulti}.
\end{enumerate}
For many coefficients, the triangular recurrence \eqref{eq:V-recursion} avoids
repeated composition and is preferable.
\end{algorithm}

\subsection{Numerical inversion}

For a numerical value $y$:
\begin{enumerate}[label=\arabic*.]
\item compute $\rho$ by $W_{-1}$;
\item compute $\mu_0$ from \eqref{eq:mu0-equation};
\item evaluate a finite paired Rademacher correction $\Rcal_K$;
\item solve
$F_0(\mu)+\log(1+\Rcal_K(\mu))=\log(y/\kappa)$ by Newton's method;
\item return $N(\mu)$, and use \eqref{eq:rounding} for the sequence inverse.
\end{enumerate}
The derivative is asymptotic to $1$ in the $\mu$ coordinate, so the iteration is
well conditioned.

\section{Numerical validation}

All computations in this section used integer partition values generated by
Euler's pentagonal recurrence, exact rational Dedekind sums, and 80-decimal
floating-point evaluation of the transcendental factors.  The tables are
intended as independent checks of normalization, sector ordering, and inverse
signs; they are not used in any proof.

\subsection{Forward expansion}

Let $H(n)$ be the elementary Hardy--Ramanujan term \eqref{eq:HR}, let
$D(n)=P_{1,+}(n)$, and let $S_K(n)$ be the paired Rademacher sum through $k=K$.
The entries below are relative errors $X(n)/p(n)-1$.

\begin{table}[htbp]
\centering
\scriptsize
\setlength{\tabcolsep}{3.8pt}
\caption{Relative errors of direct approximations.}
\label{tab:forward-errors}
\begin{tabular}{rrrrrr}
\toprule
$n$ & $H$ & $D$ & $S_2$ & $S_3$ & $S_5$\\
\midrule
10   & $ 1.4534069\!\times10^{-1}$ & $-8.8622418\!\times10^{-3}$ & $ 1.6605794\!\times10^{-3}$ & $ 3.7391400\!\times10^{-4}$ & $ 3.4721733\!\times10^{-4}$\\
50   & $ 6.5439754\!\times10^{-2}$ & $-7.3073551\!\times10^{-5}$ & $ 3.9027635\!\times10^{-6}$ & $ 2.1049640\!\times10^{-7}$ & $-1.1350427\!\times10^{-7}$\\
100  & $ 4.5713563\!\times10^{-2}$ & $-1.8219969\!\times10^{-6}$ & $ 8.6837718\!\times10^{-9}$ & $-4.9491206\!\times10^{-9}$ & $ 3.1453999\!\times10^{-10}$\\
500  & $ 2.0094903\!\times10^{-2}$ & $-2.4389452\!\times10^{-13}$ & $ 1.7536653\!\times10^{-17}$ & $-2.0016640\!\times10^{-19}$ & $-3.5933306\!\times10^{-22}$\\
1000 & $ 1.4152439\!\times10^{-2}$ & $-1.6997051\!\times10^{-18}$ & $ 1.2573854\!\times10^{-24}$ & $-3.4743827\!\times10^{-27}$ & $ 3.9452980\!\times10^{-30}$\\
\bottomrule
\end{tabular}
\end{table}

The contrast between $H$ and $D$ confirms that the full algebraic
half-power series belongs to the single shifted dominant sector.  The remaining
error is exponentially small and alternates with the arithmetic amplitudes.

\subsection{Inverse expansion and phase locking}

For $y=p(n)$, define four approximations:
\begin{enumerate}[label=(\roman*)]
\item $N_\rho=1/24+3\rho^2/(2\pi^2)$;
\item $N_0$, the exact dominant-core inverse;
\item $N_{\mathrm{lin},10}:=N_0+(3\mu_0/\pi^2)v_{\mathrm{lin}}$, where
      $v_{\mathrm{lin}}=-\Rcal_{10}(\mu_0)/F_0'(\mu_0)$; this retains only
      the first-order inverse-index coefficient;
\item $N_{20}$, the Newton inverse of the paired interpolation truncated at
      $k=20$.
\end{enumerate}
The table lists approximation minus $n$.

\begin{table}[htbp]
\centering
\small
\caption{Errors in the inverse at exact partition values $y=p(n)$.}
\label{tab:inverse-errors}
\begin{tabular}{r@{\quad}rrrr}
\toprule
$n$ & $N_\rho-n$ & $N_0-n$ & $N_{\mathrm{lin},10}-n$ & $N_{20}-n$\\
\midrule
10   & $-3.9910002\!\times10^{-1}$ & $ 2.8447205\!\times10^{-2}$ & $-8.0725131\!\times10^{-4}$ & $-1.0235155\!\times10^{-4}$\\
50   & $-3.5044252\!\times10^{-1}$ & $ 4.5102431\!\times10^{-4}$ & $-9.9453881\!\times10^{-8}$ & $ 1.9795735\!\times10^{-8}$\\
100  & $-3.3599875\!\times10^{-1}$ & $ 1.5378046\!\times10^{-5}$ & $-5.2359563\!\times10^{-11}$ & $ 7.2911158\!\times10^{-11}$\\
500  & $-3.1767515\!\times10^{-1}$ & $ 4.4042295\!\times10^{-12}$ & $ 1.1396547\!\times10^{-23}$ & $ 1.6512140\!\times10^{-23}$\\
1000 & $-3.1356078\!\times10^{-1}$ & $ 4.2959998\!\times10^{-17}$ & $-3.8971591\!\times10^{-31}$ & $ 3.37\!\times10^{-35}$\\
\bottomrule
\end{tabular}
\end{table}

The first column tends slowly toward $-3/\pi^2$, exactly as predicted by
\eqref{eq:phase-shift}.  The dominant-core correction removes that $O(1)$
phase error, after which the residual is of Rademacher size.  Linearized
nonperturbative inversion already captures many additional digits.

\subsection{Parity check}

At $n=49,50,51$, the dominant-core errors are respectively
\[
 -5.2643134\times10^{-4},\qquad
  4.5102431\times10^{-4},\qquad
 -4.0920475\times10^{-4}.
\]
Their alternating sign is explained by
$\Acal_2(n)=(-1)^n$.  After applying only the $k=2,+$ correction
\eqref{eq:Delta2}, the errors become
\[
 -1.0889617\times10^{-5},\qquad
 -2.4087145\times10^{-5},\qquad
  2.9028845\times10^{-5},
\]
and adding the remaining linear sectors through $k=10$ reduces them to
$-2.24\times10^{-7}$, $-9.95\times10^{-8}$, and $-5.81\times10^{-8}$.
This verifies both the sign and the phase of the leading inverse chirp.

\section{Interpretation and further deductions}

\subsection{A two-level transseries geometry}

The direct expansion has a nested geometry.  At the outer level, the dominant
sector is multiplied by $1+\Rcal$.  At the inner level, the coefficient of the
condensed factor $\e^{-\mu}$ in \eqref{eq:R-exact} contains
$\e^{\pm\mu/k}$ and is resolved by increasing $k$.  This is analogous to a
transasymptotic expansion in which an infinite family of saddles condenses at a
limiting action.  The pairing constraint is part of that geometry.

The inverse inherits the same structure.  Product sectors have actions
\begin{equation}\label{eq:multi-action}
 \alpha(\multi m)=\sum_{j\in\Jcal_K}m_j\alpha_j.
\end{equation}
Below action $1$, only the primary positive sectors occur: any product of two
positive actions is at least $1$.  At action $1$, nonlinear products such as
$(k=2,+)^2$ appear at the same limiting level approached by the infinitely many
primary positive sectors.  This is where the paired completion becomes
essential.

\subsection{Arithmetic phase as a coefficient algebra}

After phase locking, the coefficient algebra is generated by rational
functions of $\mu_0$ and the finite trigonometric polynomials
\begin{equation}\label{eq:chirp-generator}
 \Acal_k(N_0)=\sum_{(h,k)=1}
 \cos\!\left(\pi s(h,k)-\frac{2\pi h}{k}
 \left(\frac1{24}+\frac{3\mu_0^2}{2\pi^2}\right)\right).
\end{equation}
Differentiation with respect to $\mu_0$ preserves this algebra after adjoining
polynomial factors in $\mu_0$.  Therefore every Taylor coefficient in
\cref{thm:Vmulti} remains explicit.  In terms of $y$, the phases are quadratic
in $\log y$ to leading order.

\subsection{Why the inverse is unusually accurate}

The leading partition value grows like $\e^{c\sqrt n}$, so a relative error
$\varepsilon$ in $p(n)$ induces an index error of order $\sqrt n\,\varepsilon$.
The first omitted direct sector after the dominant core has relative size
$\e^{-c\sqrt n/2}$.  Hence the dominant-core inverse at $y=p(n)$ differs from
$n$ by approximately
\[
 O\!\left(\sqrt n\,\e^{-c\sqrt n/2}\right),
\]
which rapidly becomes much smaller than one.  This explains why rounding a
carefully shifted inverse can recover the exact index far earlier than the
slow relative convergence of the elementary Hardy--Ramanujan approximation
might suggest.

\subsection{Generalization template}

The same method applies to a broad class of coefficient sequences with a
Rademacher- or saddle-sector representation
\begin{equation}\label{eq:general-template}
 a(n)=K\e^{\Phi(n)}\Phi(n)^{-\gamma}
 \left(A_0(\Phi(n))+\sum_jA_j(n,\Phi(n))\e^{-\alpha_j\Phi(n)}\right).
\end{equation}
The reusable steps are:
\begin{enumerate}[label=\arabic*.]
\item choose a phase variable that makes each saddle prefactor elementary;
\item derive all algebraic coefficients with Bell polynomials;
\item invert the dominant exponential-power carrier with Lambert $W$ or its
      generalized analogue;
\item resum every $O(1)$ displacement in any periodic coefficient phase;
\item apply the multi-index Lagrange formula to the exponentially small
      sectors;
\item retain a paired or integral completion whenever saddle actions
      accumulate.
\end{enumerate}

\section{Conclusion}

The partition numbers possess a particularly transparent complete asymptotic
structure once the modular shift $n-1/24$ is used.  Every fixed Rademacher
sector is exactly a two-term exponential expression, and its entire
$n^{-1/2}$ coefficient sequence is available in the finite formula
\eqref{eq:C-finite}.  The ordinary Hardy--Ramanujan Poincare series is only the
algebraic expansion of the dominant positive sector; the rest of Rademacher's
convergent series is its nonperturbative completion.

For inversion, the exact Lambert carrier \eqref{eq:rho-W}, the all-order
algebraic correction \eqref{eq:a-lagrange}--\eqref{eq:a-recurrence}, and the
multi-index nonperturbative formula \eqref{eq:Vmulti-compact} provide a complete
coefficient calculus.  The arithmetic amplitudes force a phase-locked
organization based at $N_0(y)$ rather than the raw Lambert index.  The first
correction is the parity chirp \eqref{eq:leading-inverse-chirp}; subsequent
primary sectors form a fractional-power staircase accumulating at $y^{-1}$.
The paired Rademacher tail completes the expansion at that accumulation level.

The resulting framework is both symbolic and analytic: coefficients are
finite combinatorial expressions, while residual transfer and the explicit
paired-tail estimate convert truncations into controlled approximations.  It
also gives a precise bridge from a smooth inverse to the generalized inverse of
the discrete sequence through the exact ceiling relation
\eqref{eq:rounding}.

\appendix

\section{Residue proof of the finite sector coefficient formula}
\label{app:finite-sum-proof}

This appendix supplies the promised detailed proof of
\cref{thm:C-finite}.  Start from the generating function
\eqref{eq:C-generating}.  Put
\begin{equation}\label{eq:appendix-u-b}
 u=-\frac{t}{4\sqrt6},\qquad b=\frac{\sigma\pi}{6k}.
\end{equation}
Because $\sigma c/k=2\sqrt6\,b$, one obtains
\begin{equation}\label{eq:G-u}
 G(-4\sqrt6u):=\sum_{r\geq0}C_{k,\sigma;r}(-4\sqrt6u)^r
 =\frac{\e^{bT(u)}}{1-4u^2}
 \left(1+\frac{2u}{b\sqrt{1-4u^2}}\right),
\end{equation}
where
\begin{equation}\label{eq:T-catalan}
 T(u):=\frac{1-\sqrt{1-4u^2}}{2u}
\end{equation}
satisfies
\begin{equation}\label{eq:T-functional}
 T=u(1+T^2),\qquad
 u=\frac{T}{1+T^2},\qquad
 \sqrt{1-4u^2}=\frac{1-T^2}{1+T^2}.
\end{equation}
Differentiating the last expression for $u$ gives
\begin{equation}\label{eq:du-dT}
 \dd u=\frac{1-T^2}{(1+T^2)^2}\dd T.
\end{equation}
Substitution into \eqref{eq:G-u} yields the exact differential identity
\begin{align}
 G(-4\sqrt6u)\dd u
 &=\e^{bT}
 \left(\frac1{1-T^2}+\frac{2T}{b(1-T^2)^2}\right)\dd T
 \nonumber\\
 &=\frac1b\dd\left(\frac{\e^{bT}}{1-T^2}\right).
 \label{eq:G-differential}
\end{align}
Let
\begin{equation}\label{eq:Dr}
 D_r:=C_{k,\sigma;r}(-4\sqrt6)^r=[u^r]G(-4\sqrt6u).
\end{equation}
Residue transfer gives
\begin{align}
 D_r
 &=\Res_{u=0}G(-4\sqrt6u)u^{-r-1}\dd u\nonumber\\
 &=\frac1b\Res_{T=0}
 T^{-r-1}(1+T^2)^{r+1}
 \dd\left(\frac{\e^{bT}}{1-T^2}\right).
 \label{eq:Dr-residue-1}
\end{align}
The residue of a derivative is zero, so integration by parts gives
\begin{align}
 D_r
 &=-\frac1b\Res_{T=0}\frac{\e^{bT}}{1-T^2}
 \dd\left(T^{-r-1}(1+T^2)^{r+1}\right)\nonumber\\
 &=\frac{r+1}{b}\Res_{T=0}
 \e^{bT}(1+T^2)^rT^{-r-2}\dd T\nonumber\\
 &=\frac{r+1}{b}[T^{r+1}]\e^{bT}(1+T^2)^r.
 \label{eq:Dr-residue-2}
\end{align}
Expanding the exponential and binomial gives
\begin{equation}\label{eq:Dr-binomial}
 D_r=\frac{r+1}{b}
 \sum_{j=0}^{\floor{(r+1)/2}}
 \binom rj\frac{b^{r+1-2j}}{(r+1-2j)!}.
\end{equation}
Finally,
\begin{equation}\label{eq:binomial-identity}
 (r+1)\binom rj=(r+1-j)\binom{r+1}{j},
\end{equation}
and division by $(-4\sqrt6)^r$ proves \eqref{eq:C-finite}.

\section{Additional algebraic inverse coefficients}

For reference, the next coefficients after those displayed in
\eqref{eq:a-table-inline} are
\begin{equation}\label{eq:a-more}
 a_{11}=\frac{29813143}{18480},\qquad
 a_{12}=\frac{496569589}{356400}.
\end{equation}
The corresponding continuation of \eqref{eq:N0-rho} is
\begin{equation}\label{eq:N0-more}
 -\frac{101329}{420\pi^2\rho^7}
 -\frac{228677}{3360\pi^2\rho^8}+O(\rho^{-9}).
\end{equation}

The recurrence \eqref{eq:a-recurrence} can be checked directly from the first
few ordinary Bell polynomials:
\begin{align*}
 \Bhat_{n,1}&=a_n,\\
 \Bhat_{n,2}&=\sum_{j=1}^{n-1}a_ja_{n-j},\\
 \Bhat_{n,3}&=\sum_{i+j+k=n}a_ia_ja_k,
\end{align*}
with all indices positive.  This representation is useful for exact rational
computation.

\section{Formula sheet}

For convenient reuse, the central formulas are collected here.

\begin{longtable}{p{0.31\textwidth}p{0.63\textwidth}}
\toprule
Object & Formula\\
\midrule
\endfirsthead
\toprule
Object & Formula\\
\midrule
\endhead
Shift and phase & $x=n-1/24$, $\mu=c\sqrt x$, $c=\pi\sqrt{2/3}$\\[1mm]
Exact sector & $P_{k,\sigma}=A_k(4\sqrt{3k}\,x)^{-1}
\e^{\sigma\mu/k}(1-\sigma k/\mu)$\\[1mm]
Action & $\alpha_{k,\sigma}=1-\sigma/k$\\[1mm]
Paired kernel & $g(z)=\cosh z-\sinh z/z$\\[1mm]
Dominant factor & $D(\mu)=\kappa\e^\mu\mu^{-2}(1-\mu^{-1})$,
$\kappa=\pi^2/(6\sqrt3)$\\[1mm]
Direct coefficient & Equation \eqref{eq:C-finite}\\[1mm]
Bell coefficient & Equation \eqref{eq:C-bell}\\[1mm]
Lambert carrier & $\rho=-2W_{-1}(-\frac12\sqrt{\kappa/y})$\\[1mm]
Carrier polynomial & Equation \eqref{eq:P-stirling}\\[1mm]
Algebraic core & $\mu_0=\rho+U(\rho^{-1})$, with \eqref{eq:U-fixed}\\[1mm]
Core coefficient & Equations \eqref{eq:a-lagrange} and \eqref{eq:a-recurrence}\\[1mm]
Full one-stage inverse & Equation \eqref{eq:one-stage-inverse}\\[1mm]
Multi-sector coefficient & Equation \eqref{eq:Vmulti-compact}\\[1mm]
Inverse-index coefficient & Equation \eqref{eq:Nmulti}\\[1mm]
Leading chirp & Equation \eqref{eq:leading-inverse-chirp}\\[1mm]
Discrete inverse & $p_{\Z}^{\langle-1\rangle}(y)=\ceil{\Ncal(y)}$\\
\bottomrule
\end{longtable}

\section{Notes on reproducibility}

The symbolic identities can be verified with any system that supports formal
power series.
\begin{enumerate}[label=\arabic*.]
\item Expand \eqref{eq:C-generating} and compare with
      \eqref{eq:C-finite}; equality through order $r$ uses exact arithmetic in
      $b$.
\item Substitute $U=\sum_{m=1}^Ma_ms^m$ into \eqref{eq:U-fixed}; the residual
      begins at $s^{M+1}$.
\item Substitute \eqref{eq:rho-log-expansion} with \eqref{eq:P-stirling} into
      $\rho-2\log\rho-L$; the residual begins at the requested power of $L^{-1}$.
\item For a finite sector set, substitute coefficients from
      \eqref{eq:Vmulti-compact} into \eqref{eq:v-equation-expanded}; every
      prescribed multi-index cancels identically.
\end{enumerate}
Numerical partition values may be generated without external tables by Euler's
pentagonal recurrence
\begin{equation}\label{eq:pentagonal-recurrence}
 p(n)=\sum_{m\geq1}(-1)^{m-1}
 \left[p\!\left(n-\frac{m(3m-1)}2\right)
      +p\!\left(n-\frac{m(3m+1)}2\right)\right],
\end{equation}
where $p(0)=1$ and $p(r)=0$ for $r<0$.

\begingroup
\small
\setlength{\itemsep}{0.15em}
\begin{thebibliography}{99}

\bibitem{HardyRamanujan1918}
G.~H. Hardy and S. Ramanujan,
\newblock Asymptotic formulae in combinatory analysis,
\newblock \emph{Proceedings of the London Mathematical Society} (2) 17
(1918), 75--115.

\bibitem{Rademacher1937PNAS}
H. Rademacher,
\newblock A convergent series for the partition function $p(n)$,
\newblock \emph{Proceedings of the National Academy of Sciences USA} 23
(1937), 78--84.

\bibitem{Rademacher1938}
H. Rademacher,
\newblock On the partition function $p(n)$,
\newblock \emph{Proceedings of the London Mathematical Society} (2) 43
(1938), 241--254.

\bibitem{Rademacher1943}
H. Rademacher,
\newblock On the expansion of the partition function in a series,
\newblock \emph{Annals of Mathematics} 44 (1943), 416--422.

\bibitem{Lehmer1939}
D.~H. Lehmer,
\newblock On the remainders and convergence of the series for the partition
function,
\newblock \emph{Transactions of the American Mathematical Society} 46
(1939), 362--373.

\bibitem{OSullivan2023}
C. O'Sullivan,
\newblock Detailed asymptotic expansions for partitions into powers,
\newblock arXiv:2205.13468v2, 2023,
\newblock \url{https://arxiv.org/abs/2205.13468}.

\bibitem{BanerjeeEtAl2022}
K. Banerjee, P. Paule, C.-S. Radu, and C. Schneider,
\newblock Error bounds for the asymptotic expansion of the partition function,
\newblock arXiv:2209.07887, 2022,
\newblock \url{https://arxiv.org/abs/2209.07887}.

\bibitem{KongTeo2023}
Z.-Y. Kong and L.-P. Teo,
\newblock Rademacher's formula for the partition function,
\newblock arXiv:2302.03835, 2023,
\newblock \url{https://arxiv.org/abs/2302.03835}.

\bibitem{CorlessEtAl1996}
R.~M. Corless, G.~H. Gonnet, D.~E.~G. Hare, D.~J. Jeffrey, and D.~E. Knuth,
\newblock On the Lambert $W$ function,
\newblock \emph{Advances in Computational Mathematics} 5 (1996), 329--359.

\bibitem{JeffreyEtAl1995}
D.~J. Jeffrey, R.~M. Corless, D.~E.~G. Hare, and D.~E. Knuth,
\newblock On the inversion of $y^\alpha e^y$ in terms of associated Stirling
numbers,
\newblock arXiv:math/9512230, 1995.

\bibitem{CombinatorialCoefficientCalculus}
V. Reshetnikov,
\newblock \emph{Combinatorial Coefficient Calculus},
\newblock ProveIt/Fabius Function research corpus, current repository version,
\newblock \url{https://github.com/VladimirReshetnikov/ProveIt/tree/main/Analysis/FabiusFunction/docs/semi-formalized-research-frontiers/drafts/combinatorial-coefficient-calculus/Combinatorial_Coefficient_Calculus}.

\bibitem{PolynomialLogTransseries}
V. Reshetnikov,
\newblock \emph{Polynomial--Logarithmic Transseries},
\newblock ProveIt/Fabius Function research corpus, current repository version,
\newblock \url{https://github.com/VladimirReshetnikov/ProveIt/tree/main/Analysis/FabiusFunction/docs/semi-formalized-research-frontiers/drafts/series-and-transseries}.

\bibitem{OEISA000041}
OEIS Foundation Inc.,
\newblock The On-Line Encyclopedia of Integer Sequences, A000041,
\newblock \url{https://oeis.org/A000041}.

\bibitem{EberlPaulson2026}
M. Eberl and L.~C. Paulson,
\newblock Rademacher's series for the partition function,
\newblock Archive of Formal Proofs, 3 July 2026,
\newblock \url{https://isa-afp.org/entries/Rademacher_Series.html}.

\end{thebibliography}
\endgroup

\end{document}
